% ===== PRÉAMBULE CANONIQUE — à coller VERBATIM en tête de CHAQUE .tex =====
\documentclass[11pt,a4paper]{article}
\usepackage[utf8]{inputenc}\usepackage[T1]{fontenc}\usepackage{lmodern}
\usepackage[french]{babel}
\usepackage{amsmath,amssymb,mathtools}\usepackage{icomma}
\usepackage[version=4]{mhchem}          % \ce{...} : équations & formules
\usepackage{siunitx}\sisetup{locale=FR} % \qty{}{}, \si{}, \ang{}
\usepackage{chemfig}                    % \chemfig{...} : structures développées
\usepackage{xcolor}\usepackage{colortbl}
\usepackage{tikz}\usepackage{pgfplots}\pgfplotsset{compat=1.18}
\usepackage[most]{tcolorbox}\usepackage{enumitem}\usepackage{array,booktabs}
\usepackage[a4paper,margin=2.2cm]{geometry}
\usetikzlibrary{arrows.meta,calc,angles,quotes,positioning,patterns,backgrounds,shapes.geometric}
\definecolor{primary}{RGB}{222,49,99}\definecolor{primarydark}{RGB}{150,22,55}
\definecolor{defcol}{RGB}{0,114,178}\definecolor{methcol}{RGB}{52,92,148}
\definecolor{excol}{RGB}{0,135,90}\definecolor{attncol}{RGB}{200,40,55}
\tcbset{boxbase/.style={enhanced,breakable,boxrule=0.9pt,arc=2.5pt,
  left=3mm,right=3mm,top=2.3mm,bottom=2.3mm,fonttitle=\bfseries,coltitle=white}}
% --- boîtes TUTORIEL (noms EXACTS, ne pas renommer) ---
\newtcolorbox{cle}[1][]{boxbase,colback=defcol!6,colframe=defcol,
  title={\textsc{À retenir}\ifx\relax#1\relax\relax\else\textmd{ : \itshape #1}\fi}}
\newtcolorbox{methode}[1][]{boxbase,colback=methcol!7,colframe=methcol,sharp corners,
  title={\textsc{Méthode}\ifx\relax#1\relax\relax\else\textmd{ --- \itshape #1}\fi}}
\newtcolorbox{exemplebox}[1][]{boxbase,colback=excol!5,colframe=excol,
  title={\textsc{Exemple}\ifx\relax#1\relax\relax\else\textmd{ : \itshape #1}\fi}}
\newtcolorbox{piege}[1][]{boxbase,colback=attncol!6,colframe=attncol,
  title={\textsc{Piège}\ifx\relax#1\relax\relax\else\textmd{ : \itshape #1}\fi}}
\newtcolorbox{remarque}[1][]{boxbase,colback=black!4,colframe=black!45,
  title={\textsc{Remarque}\ifx\relax#1\relax\relax\else\textmd{ : \itshape #1}\fi}}
% --- boîtes EXERCICE / CORRIGÉ (noms EXACTS, auto-comptées) ---
\newtcolorbox[auto counter]{exo}[1][]{boxbase,colback=primary!4,colframe=primary,
  title={\textsc{Exercice}~\thetcbcounter\ifx\relax#1\relax\relax\else\textmd{ --- \itshape #1}\fi}}
\newtcolorbox[auto counter]{corrige}[1][]{boxbase,colback=excol!4,colframe=excol,
  title={\textsc{Corrigé}~\thetcbcounter\ifx\relax#1\relax\relax\else\textmd{ --- \itshape #1}\fi}}
\newcommand{\R}{\mathbb{R}}\newcommand{\N}{\mathbb{N}}\newcommand{\Z}{\mathbb{Z}}\newcommand{\ds}{\displaystyle}
% =========================================================================
\begin{document}
\begin{corrige}[chaleur sensible $Q=mc\Delta T$]
\begin{enumerate}[label=\alph*)]
  \item $Q = 2{,}0\times 4{,}185\times(30-20) = 2{,}0\times 4{,}185\times 10 = \qty{83.7}{\kilo\joule}$.
  \item $Q = 0{,}50\times 4{,}185\times(100-15) = 0{,}50\times 4{,}185\times 85 = \qty{177.9}{\kilo\joule}$.
\end{enumerate}
\end{corrige}

\begin{corrige}[chaleur latente $Q=mL$]
\begin{enumerate}[label=\alph*)]
  \item $Q = m\,L_{\text{fusion}} = 0{,}20\times 334 = \qty{66.8}{\kilo\joule}$.
  \item $Q = m\,L_{\text{vaporisation}} = 0{,}10\times 2257 = \qty{225.7}{\kilo\joule}$.
\end{enumerate}
\end{corrige}

\begin{corrige}[chaleur reçue ou cédée ?]
\begin{enumerate}[label=\alph*)]
  \item $\Delta T = 60-20 = +\qty{40}{\celsius}>0$ : le corps \emph{reçoit} de la chaleur, $Q>0$.
  \item $\Delta T = 25-80 = -\qty{55}{\celsius}<0$ : le corps \emph{cède} de la chaleur, $Q<0$.
  \item enceinte adiabatique : $Q_1 + Q_2 = 0$, donc $Q_1 = -Q_2$ (ce que l'un cède, l'autre le reçoit).
\end{enumerate}
\end{corrige}

\begin{corrige}[conversions indispensables]
\begin{enumerate}[label=\alph*)]
  \item $\qty{4.185}{\kilo\joule} = 4{,}185\times 1000 = \qty{4185}{\joule}$.
  \item $\qty{250}{\gram} = \dfrac{250}{1000} = \qty{0.250}{\kilogram}$.
  \item une \emph{variation} de $\qty{25}{\celsius}$ vaut $\qty{25}{\kelvin}$ (identique pour un écart).
\end{enumerate}
\end{corrige}

\begin{corrige}[comparer deux dépenses d'énergie]
\begin{enumerate}[label=\alph*)]
  \item fondre : $Q = m\,L_{\text{fusion}} = 0{,}10\times 334 = \qty{33.4}{\kilo\joule}$.
  \item chauffer : $Q = m\,c_{\text{eau}}\,\Delta T = 0{,}10\times 4{,}185\times 50 = \qty{20.9}{\kilo\joule}$.
  \item la \textbf{fusion} coûte plus ($\qty{33.4}{\kilo\joule}>\qty{20.9}{\kilo\joule}$) : les changements d'état sont très gourmands en énergie.
\end{enumerate}
\end{corrige}

\end{document}
